\documentclass{amsart}
\usepackage{amsmath,amssymb,amsthm,hyperref}

\newtheorem{theorem}{Theorem}
\newtheorem{lemma}{Lemma}
\newtheorem{corollary}{Corollary}
\theoremstyle{remark}
\newtheorem*{remark}{Remark}

\begin{document}

\title{An independence bound for planar graphs without the Four Color Theorem}
\maketitle

\begin{center}\textbf{Honest status of the target claim}\end{center}

The problem asks to prove that every planar graph on $n$ vertices has an
independent set of size at least $n/4$, \emph{without} invoking the Four Color
Theorem (4CT). I must report at the outset that the bound $n/4$, in this
unconditional form, is \textbf{not known to be provable without the 4CT}: the
best published lower bound on the independence ratio of planar graphs that does
not use the 4CT is $3/13$ (Cranston--Rabern, 2016), and $3/13 = 0.2307\ldots <
1/4$. The bound $n/4$ is tight (the complete graph $K_4$ has $n=4$,
$\alpha=1=n/4$, and the icosahedron has $n=12$, $\alpha=3=n/4$), and obtaining it
unconditionally is widely regarded as comparable in difficulty to the 4CT
itself. In the interest of mathematical honesty, I prove below, completely and
from first principles, the strongest clean bound that \emph{is} attainable by an
elementary self-contained argument that does not use the 4CT, namely
\[
   \alpha(G) \;\ge\; \frac{n}{5},
\]
and I delineate precisely (Section~\ref{sec:gap}) where the elementary method
stops short of $n/4$ and why that is the genuine obstruction.

\bigskip

Throughout, $G=(V,E)$ is a finite simple graph with $n=|V|$ vertices and $m=|E|$
edges, $\alpha(G)$ denotes the maximum size of an independent set, and
$\chi(G)$ denotes the chromatic number. For a vertex $v$, $N(v)$ is its set of
neighbours, $\deg(v)=|N(v)|$, and $N[v]=N(v)\cup\{v\}$. We fix a plane embedding
of the planar graph $G$ when one is needed.

\section{Pigeonhole reduction from colouring to independence}

\begin{lemma}\label{lem:pigeon}
If a graph $G$ on $n$ vertices satisfies $\chi(G)\le k$, then
$\alpha(G)\ge n/k$.
\end{lemma}

\begin{proof}
By definition of $\chi(G)\le k$ there is a proper colouring
$c\colon V\to\{1,\dots,k\}$, i.e.\ $c(u)\ne c(v)$ whenever $uv\in E$. For each
colour $i\in\{1,\dots,k\}$ let $C_i=c^{-1}(i)=\{v\in V: c(v)=i\}$. The sets
$C_1,\dots,C_k$ partition $V$, so $\sum_{i=1}^{k}|C_i|=n$. By the pigeonhole
principle there is an index $j$ with $|C_j|\ge n/k$. Each $C_i$ is independent:
if $u,v\in C_i$ then $c(u)=c(v)=i$, and properness forces $uv\notin E$. Hence
$C_j$ is an independent set with $|C_j|\ge n/k$, giving $\alpha(G)\ge n/k$.
\end{proof}

Thus it suffices to prove $\chi(G)\le 5$ for every planar graph $G$. (The 4CT is
the statement $\chi(G)\le 4$; we deliberately avoid it and prove the weaker,
elementary Five Color Theorem, which yields the $n/5$ bound.)

\section{Euler's formula and a low-degree vertex}

\begin{lemma}\label{lem:edges}
If $G$ is a simple planar graph with $n\ge 3$ vertices and $m$ edges, then
$m\le 3n-6$.
\end{lemma}

\begin{proof}
It suffices to prove the inequality when $G$ is connected: if $G$ has components
of sizes $n_1,\dots,n_t$ summing to $n$, with $t\ge2$, one may add an edge
joining two distinct components without destroying planarity and without creating
a multi-edge, strictly increasing $m$ while keeping $n$ fixed; iterating, a graph
with the maximum number of edges for a given $n$ is connected, so it suffices to
bound $m$ for connected $G$. Assume then that $G$ is connected and fix a plane
embedding. Let $f$ denote the number of faces of this embedding (including the
unbounded face). Euler's formula for connected plane graphs states
\begin{equation}\label{eq:euler}
   n - m + f = 2 .
\end{equation}
If $G$ has no cycle then, being connected, it is a tree, so $m=n-1$ and (for
$n\ge3$) $m=n-1\le 3n-6$, as required. Assume now that $G$ contains a cycle. We
bound $f$. Counting incidences between faces and edge-sides: each edge borders at
most two faces, so it contributes at most $2$ to the sum of boundary-walk lengths
over all faces; hence
\[
   \sum_{\text{faces } \phi}\operatorname{len}(\partial\phi)\;\le\;2m .
\]
On the other hand, since $G$ is simple and contains a cycle, every face boundary
walk has length at least $3$ (a boundary walk of length $1$ requires a loop and
of length $2$ requires two parallel edges, neither of which occurs in a simple
graph). Therefore
\[
   3f \;\le\; \sum_{\text{faces } \phi}\operatorname{len}(\partial\phi)\;\le\;2m,
   \qquad\text{so}\qquad f\le \tfrac{2m}{3}.
\]
Substituting into \eqref{eq:euler},
\[
   2 = n-m+f \le n - m + \tfrac{2m}{3} = n - \tfrac{m}{3},
\]
which rearranges to $m\le 3n-6$.
\end{proof}

\begin{lemma}\label{lem:mindeg}
Every simple planar graph $G$ has a vertex of degree at most $5$.
\end{lemma}

\begin{proof}
If $n\le2$ this is immediate ($\deg(v)\le n-1\le1\le5$). Assume $n\ge3$ and
suppose, for contradiction, that every vertex has degree at least $6$. Then by the
handshake identity $2m=\sum_{v}\deg(v)\ge 6n$, so $m\ge 3n$. This contradicts
Lemma~\ref{lem:edges}, which gives $m\le 3n-6<3n$. Hence some vertex has degree
$\le5$.
\end{proof}

\section{The Five Color Theorem}

\begin{theorem}[Five Color Theorem]\label{thm:five}
Every planar graph $G$ satisfies $\chi(G)\le 5$.
\end{theorem}

\begin{proof}
We argue by induction on $n=|V(G)|$. If $n\le5$, assign distinct colours from
$\{1,\dots,5\}$ to the at most five vertices; this is a proper $5$-colouring.

\emph{Inductive step.} Let $n\ge6$ and assume every planar graph on fewer than
$n$ vertices is $5$-colourable. Fix a plane embedding of $G$. By
Lemma~\ref{lem:mindeg}, $G$ has a vertex $v$ with $\deg(v)\le5$. The graph $G-v$
is planar (deleting a vertex preserves planarity) and has $n-1$ vertices, so by
induction it has a proper $5$-colouring $c\colon V(G)\setminus\{v\}\to
\{1,\dots,5\}$. We extend $c$ to $v$.

\medskip
\noindent\emph{Case 1: $\deg(v)\le 4$, or the neighbours of $v$ use at most $4$
distinct colours.} Then at least one colour, say $a$, appears on no neighbour of
$v$. Set $c(v)=a$. Since $a$ differs from the colour of every neighbour of $v$
and $c$ was proper on $G-v$, the extension is a proper $5$-colouring of $G$.

\medskip
\noindent\emph{Case 2: $\deg(v)=5$ and the five neighbours of $v$ receive five
\emph{distinct} colours.} Let the neighbours, listed in the cyclic order in which
the edges $vv_i$ leave $v$ in the plane embedding, be $v_1,v_2,v_3,v_4,v_5$, and
relabel colours so that $c(v_i)=i$ for $i=1,\dots,5$.

Consider colours $1$ and $3$. Let $H_{1,3}$ be the subgraph of $G-v$ induced by
all vertices coloured $1$ or $3$, and let $K$ be the connected component of
$H_{1,3}$ containing $v_1$.

\emph{Subcase 2a: $v_3\notin K$.} Recolour, \emph{within $K$ only}, by swapping
colours $1$ and $3$ (each vertex of $K$ coloured $1$ becomes $3$ and vice versa).
Call the new colouring $c'$. We verify $c'$ is a proper colouring of $G-v$. The
only edges whose endpoints might conflict are those with exactly one endpoint in
$K$. Let $xy$ be such an edge with $x\in K$, $y\notin K$. As $x\in K\subseteq
H_{1,3}$, we have $c(x)\in\{1,3\}$. If $y$ had colour $1$ or $3$, then $y\in
H_{1,3}$ and, being adjacent to $x\in K$, $y$ would lie in the same component
$K$, contradicting $y\notin K$. Hence $c(y)\in\{2,4,5\}$. The swap gives
$c'(x)\in\{1,3\}$ and leaves $c'(y)=c(y)\in\{2,4,5\}$, so $c'(x)\ne c'(y)$. Edges
with both endpoints in $K$ keep proper colours (both endpoints swap $1\leftrightarrow3$
consistently), and edges with both endpoints outside $K$ are untouched. Hence
$c'$ is a proper $5$-colouring of $G-v$. Under $c'$, $v_1\in K$ now has colour
$3$, while $v_3\notin K$ keeps colour $3$; thus the neighbours of $v$ now use the
colours $3,2,3,4,5$, so colour $1$ appears on no neighbour of $v$. Set $c'(v)=1$
to obtain a proper $5$-colouring of $G$.

\emph{Subcase 2b: $v_3\in K$.} Then $v_1$ and $v_3$ lie in the same component of
$H_{1,3}$, so there is a path $P_{1,3}$ from $v_1$ to $v_3$ whose internal
vertices are all coloured $1$ or $3$. Adjoining the edges $vv_1$ and $vv_3$ to
$P_{1,3}$ yields a cycle $C$ drawn in the plane. By the Jordan curve theorem,
the closed curve $C$ separates the plane into a bounded interior and an unbounded
exterior. Because the embedding fixes the cyclic order $v_1,v_2,v_3,v_4,v_5$ of
edges around $v$, the neighbour $v_2$ lies in the cyclic sector strictly between
$v_1$ and $v_3$, whereas $v_4$ (and $v_5$) lie in the complementary sector;
consequently $v_2$ and $v_4$ lie in different regions determined by $C$ -- one
strictly inside $C$ and the other strictly outside. (This is exactly where
planarity is used.)

Now consider colours $2$ and $4$: let $H_{2,4}$ be the subgraph of $G-v$ induced
by vertices coloured $2$ or $4$, and $K'$ the component of $H_{2,4}$ containing
$v_2$. A path from $v_2$ to $v_4$ inside $H_{2,4}$ uses only vertices coloured
$2$ or $4$ and never uses $v$; since the curve $C$ consists solely of $v$ and
vertices coloured $1$ or $3$, such a path is disjoint from $C$ and hence cannot
cross from the interior to the exterior of $C$. As $v_2$ and $v_4$ lie on
opposite sides of $C$, no $\{2,4\}$-coloured path joins them, so $v_4\notin K'$.
We now repeat the swap of Subcase~2a with the colour pair $(2,4)$ and component
$K'$: swapping colours $2$ and $4$ on $K'$ produces a proper $5$-colouring $c''$
of $G-v$ (by the identical verification as in Subcase~2a, with $\{1,3\}$ replaced
by $\{2,4\}$) under which $v_2$ becomes colour $4$ while $v_4$ retains colour $4$.
Then colour $2$ appears on no neighbour of $v$, and setting $c''(v)=2$ gives a
proper $5$-colouring of $G$.

In every case $G$ admits a proper $5$-colouring, completing the induction.
\end{proof}

\section{The independence bound}

\begin{theorem}\label{thm:main}
Every planar graph $G$ on $n$ vertices contains an independent set of size at
least $n/5$; that is, $\alpha(G)\ge n/5$. This argument does not use the Four
Color Theorem.
\end{theorem}

\begin{proof}
By Theorem~\ref{thm:five}, $\chi(G)\le5$. Applying Lemma~\ref{lem:pigeon} with
$k=5$ gives $\alpha(G)\ge n/5$. The only colouring input is the Five Color
Theorem, whose proof uses only Euler's formula, the existence of a low-degree
vertex, and Kempe-chain recolouring; the Four Color Theorem is never invoked.
\end{proof}

\section{Where the elementary method stops short of \texorpdfstring{$n/4$}{n/4}}
\label{sec:gap}

For transparency we record exactly why the elementary toolkit yields $n/5$ but
not $n/4$.

\medskip
\noindent\textbf{(i) The colouring route caps at $n/5$.}
Lemma~\ref{lem:pigeon} turns a $k$-colouring into an independent set of size
$n/k$. Without the 4CT the best elementary colouring bound is $\chi\le5$
(Theorem~\ref{thm:five}), giving exactly $n/5$. The 4CT would give $\chi\le4$ and
hence $n/4$, but it is forbidden here.

\medskip
\noindent\textbf{(ii) The greedy/degeneracy route caps at $n/6$.}
A planar graph is $5$-degenerate (every subgraph has a vertex of degree $\le5$ by
Lemma~\ref{lem:mindeg}), which yields only $\alpha\ge n/6$ by the standard greedy
elimination.

\medskip
\noindent\textbf{(iii) A complete $n/4$ argument exists only when low-degree
vertices keep appearing.}
The bound $n/4$ is recovered cleanly for planar graphs that are
\emph{$3$-degenerate}, i.e.\ every subgraph has a vertex of degree at most $3$.
Indeed we prove by induction on $n$ that every $3$-degenerate graph $G$ satisfies
$\alpha(G)\ge n/4$. If $n\le4$ then a single vertex is an independent set and
$1\ge n/4$. Otherwise pick a vertex $v$ with $\deg(v)\le3$. The induced subgraph
$G'=G-N[v]$ has at least $n-4$ vertices and is again $3$-degenerate (every
subgraph of $G$ is), so by induction $\alpha(G')\ge (n-4)/4$. Adjoining $v$ to a
maximum independent set of $G'$ (legitimate since $v$ has no neighbours in $G'$)
gives an independent set of $G$ of size $1+\alpha(G')\ge 1+(n-4)/4=n/4$. This is a
\emph{complete, correct} proof of $\alpha\ge n/4$ for $3$-degenerate planar
graphs.

The obstruction to the general case is precisely the existence of planar graphs of
minimum degree $\ge4$ that are \emph{not} $3$-degenerate -- equivalently, graphs
in which every closed neighbourhood removal deletes at least $5$ vertices,
destroying the ``one independent vertex per $4$ deleted'' budget. The icosahedron
($5$-regular, $n=12$, $\alpha=3=n/4$) is the canonical tight obstruction: it is
$5$-regular, so the recursion above never starts.

\medskip
\noindent\textbf{(iv) State of the art.}
Closing the gap unconditionally is, to current knowledge, open. The best
published independence ratio for planar graphs avoiding the 4CT is
$3/13=0.2307\ldots<1/4$ (Cranston--Rabern, 2016), and the Albertson--Berman
conjecture -- that every planar graph has an induced forest on $\ge n/2$ vertices,
which would yield $\alpha\ge n/4$ because a forest is bipartite -- remains
unproven. Accordingly, a fully unconditional elementary proof of the exact bound
$n/4$ for \emph{all} planar graphs is not currently available. The strongest
complete results obtainable by self-contained means are: $\alpha(G)\ge n/5$ for
all planar $G$ (Theorem~\ref{thm:main}), and $\alpha(G)\ge n/4$ for all
$3$-degenerate planar $G$ (item~(iii)).

\end{document}
